% GDZ PPN236018647,native553–554before83;HKUST554–555.
\clearpage\SourceLeaf{553}\ArithmeticHead{524}{L}{TAGEBUCH MIT ERLÄUTERUNGEN.}
\NumberedItem{81}
\indent Nova theorematis \textsc{Pythagoraei} dem[onstratio].
\par\noindent\hfill Bruns[vigae, 1797] Oct. 16.

\medskip\begingroup\small
\indent Der Beweis ist im Nachlaß vorhanden:
\par\endgroup
\begin{center}
\textbf{Nova theorematis pythagoraei demonstratio.}\par\medskip
\rule{.31\linewidth}{.35pt}\par
{\small[Zettel in Ff, Kapsel 41.]}\par
\rule{.31\linewidth}{.35pt}
\end{center}
\begin{center}
\begin{tikzpicture}[x=.92cm,y=.92cm,line width=.55pt]
\coordinate (D) at (0,0);\coordinate (A) at (1.9,0);
\coordinate (B) at (4.8,0);\coordinate (F) at (9.5,0);
\coordinate (C) at (1.9,3.8);
\draw (A)--(C)--(B)--cycle;
\draw[dash pattern=on 3pt off 3pt] (D)--(A) (B)--(F) (D)--(C)--(F);
\draw[dash pattern=on 3pt off 3pt] (0,0)..controls(.01,.28)and(.04,.54)..(.07,.78);
\draw[dash pattern=on 3pt off 3pt] (9.5,0)..controls(9.51,.28)and(9.46,.62)..(9.39,.92);
\draw[dash pattern=on 3pt off 3pt] (1.48,3.52)..controls(1.76,3.76)and(2.05,3.99)..(2.34,4.16);
\node[below] at (D) {$D$};\node[below] at (A) {$A$};
\node[below] at (B) {$B$};\node[below] at (F) {$F$};
\node[above left,inner sep=2pt] at (C) {$C$};
\end{tikzpicture}
\end{center}
\begin{center}\textls[120]{Theorema.}\end{center}
\indent Si trianguli \(ABC\) angulus sit rectus, erit
\[
AB^{\mathrm{qu}}+AC^{\mathrm{qu}}=BC^{\mathrm{qu}}.
\]
\noindent\textls[120]{Demonstr[atio].}
\par\indent Producatur \(AB\) utrinque fiatque
\[
BD=BF=BC.
\]
\noindent Tum triangula \(CBD,\ CBF\) erunt isocelia et anguli
\[
BDC=BCD;\quad BCF=BFC.
\]
\noindent At
\[
BDC+DCB+BCF+CFB=2\text{ Rectis.}
\]
\noindent Quare \(DCB+BCF\) erit \(=\) Recto. Hinc Angulus \(CDA=ACF\) et triangula

\clearpage\SourceLeaf{554}\ArithmeticHead{525}{R}{1797 OCT. 16—1798 APR.}
\noindent\(ACD,\ AFC\) similia. Quare
\[
AC:AD=AF:AC
\]
\noindent et \(AC^{\mathrm{qu}}=\) Rectg. sub \(AD\) et \(AF\) i. e. sub \(BC-AB\) et \(BC+AB\). Hinc (\textsc{Eucl.} Elem. II)
\[
AC^{\mathrm{qu}}=BC^{\mathrm{qu}}-AB^{\mathrm{qu}}
\]
\noindent et
\[
AC^{\mathrm{qu}}+AB^{\mathrm{qu}}=BC^{\mathrm{qu}}\qquad\qquad QED.
\]

\NumberedItem{82}
\indent Seriei
\[
x-\frac12x^2+\frac1{12}x^3-\frac1{144}x^4\ldots
\]
\noindent summam consideravimus invenimusque eam \(=0\), si
\[
2\surd x+\frac3{16}\frac1{\surd x}-\frac{21}{1024}\frac1{\surd{.}3x}\ldots=\left(k+\frac14\right)\pi.
\]
\par\noindent\hfill Brunsv[igae, 1797] Oct. 16.

\medskip\begingroup\small
\indent Eine Erläuterung und Bestätigung dieser Aufzeichnung findet man oben S. 388—389. Hier möge nur wiederholt werden, daß in der Gleichung linker Hand im dritten Gliede statt \(\surd{.}3x\) gelesen werden muß \(\sqrt{x^3}\).
\par\noindent\hfill\textsc{Schlesinger.}
\par\endgroup
